\label{sec:poslaw}

\subsection{Question}

How does positional decoding precision scale with the number of gap genes read
jointly, and does the scaling survive a doubling of maternal \emph{bcd} dosage?
Published reference points: single gap genes localize to several percent egg
length (\%EL) \cite{dubuis2013}, while joint decoding of four gap genes reaches
$\sim$1\,\%EL \cite{petkova2019}. The intermediate regime - two and three genes -
and the dosage dependence of the \emph{scaling law itself} are unmeasured.

\subsection{Setup}

From the Liu 2013 immunofluorescence archive (Study C data) we take the two
sessions that image three gap genes simultaneously (Hb, Eve, Kr): 244 embryos at
1$\times$ \emph{bcd} dosage (line 1XA) and 142 embryos at 2$\times$ dosage (line
2XA). Position is decoded with $k$-nearest-neighbour regression ($k=5$) from
expression vectors at 50 subsampled axis positions, holding out 20\% of embryos;
every gene subset of size 1, 2 and 3 is decoded independently (2 seeds per
subset). Controls: circular position permutation within each embryo, preserving
expression autocorrelation while destroying position information.

\subsection{Result 1: error falls as $d^{-2}$ and extrapolates to the published
4-gene precision}

At 1$\times$ dosage the median absolute decoding error falls from 21.6\,\%EL
(one gene) to 6.2\,\%EL (two genes) to 2.4\,\%EL (three genes), a log-log slope
of $\alpha = 1.98 \pm$ (three-point fit). Extrapolating the same law one step to
$d=4$ predicts 1.4\,\%EL, consistent with the $\sim$1\,\%EL reported by
\citetitle{petkova2019} for joint four-gene decoding on different data with
Bayes-optimal decoding. Because our $k$NN decoder is a strict lower bound on
decoding precision, this agreement is conservative. We do not claim to beat the
published number; the contribution is the measured scaling law in the
unmeasured intermediate regime.

\subsection{Result 2 (discovery): doubling \emph{bcd} dosage shallows the
scaling exponent}

At 2$\times$ dosage the same measurement gives 21.0\,\%EL, 8.0\,\%EL and
4.2\,\%EL for one, two and three genes: a shallower exponent ($\alpha = 1.46$
vs.\ 1.98) and \emph{worse} three-gene precision (4.2 vs.\ 2.4\,\%EL), with
$d=4$ extrapolation 2.8\,\%EL. Doubling the morphogen does not add positional
information through the gap layer; it degrades multi-gene integration. This is
consistent with Study C's dosage-dependent length-scale change (DCLS) and adds a
second, independent axis on which dosage compensation fails.

We note a boundary condition: the pair-level improvement reported earlier
(Hb+Eve, 6.2 vs.\ 9.9\,\%EL across line sets) used \emph{different} line sets
(BNT vs.\ non-BNT), while the present contrast holds the three-gene session
fixed. Dosage effects on decodability are therefore line-set dependent, and the
two results do not commute.

\subsection{Controls and limits}

Structure-preserved permutation nulls sit at 24-26\,\%EL for every subset size.
Single-gene decoding (21.6\,\%EL) is barely above its null (24.4\,\%EL): on this
immunofluorescence data, positional information is almost entirely a
\emph{multi-gene} property. Limits: only two three-gene sessions exist in the
archive; the exponent is a three-point fit; the $d=4$ point is a one-step
extrapolation; $k$NN is not the optimal decoder.
